%#! latex main.tex
%\addcontentsline{toc}{chapter}{Introduction.}
%\chapter*{Introduction.}
\markboth{INTRODUCTION}{}

\begin{center}
 {\large\bf Introduction.}
\end{center}

%\addtocontents{toc}{\protect\noindent Introduction.\hfill 1}
\addcontentsline{toc}{subsection}{Introduction.}
\subsection*{Purpose of this book.}

This book stems from lectures given at the Universities of T\={o}hoku
and Hokkaid\={o}.  

Its main purpose is to introduce the reader to $p$-adic analytic
geometry and to the theory of $p$-adic analytic functions and
differential equations, by focusing on the theme of {\it period
mappings}. Of course, this approach is not meant to replace more
systematic expositions of $p$-adic analysis or geometry found in a
number of good treatises, but to complement them.

 It is our general rule, in this book, to follow as closely as possible
the complex theory, and to go back and forth between the complex and the
$p$-adic worlds. We hope that this will make the text of interest both
to some complex geometers and to some arithmetic geometers.

In the course of Chapter I, this approach will eventually become a
kaleidoscope of half-correspondences and broken echoes. We hope that the
reader will then have gained enough hindsight and wariness about these
analogies, and will enjoy seeing unity being restored at a deeper level
in Chapters II and III.

\bigskip

We have chosen the theme of period mappings because of its
central role in the nineteen-century mathematics as a fertile place of
interaction between differential equations, group theory, algebraic and
differential geometry, topology (and even number theory).  In fact, it
was a guiding thread in the early harmonious development of these
branches of mathematics, from Gauss and Riemann to Klein and Poincar\'e
(\cf [{\bf Gray86}]).

\bigskip

The origins lie in Gauss' largely unpublished work on elliptic
functions on one hand (rediscovered and extended by Abel and by Jacobi),
and on the hypergeometric differential equation in the complex domain on
the other hand. Gauss knew the connection between the two topics: the
inverse of the indefinite integral
$\int\frac{dx}{\sqrt{(1-x^2)(1-k^2x^2)}}$ is a single-valued function
with two independent periods $\omega_1$ and $\omega_2$ which are
solutions of the 
hypergeometric differential equation with parameters
$(\frac{1}{2},\frac{1}{2};1)$ in the variable $k^2$.

Riemann's point of view of the `Riemann surface' of a multivalued
complex function has given a geometric framework for all of complex
analysis. He applied this idea with equal success to Jacobi's inversion
problem for more general indefinite algebraic integrals on one hand; and
to the elucidation of the paradoxical polymorphism of hypergeometric
functions which had puzzled Gauss and Kummer by introducing the concept
of monodromy, on the other hand. He studied in detail the monodromy of
the multivalued map $ \; k^2\mapsto \tau ={\omega_2/\omega_1}\;$, the
first and basic example of a `period mapping' , whose inverse is
single-valued. He also showed (and Schwarz rediscovered) that under
certain conditions, the quotient of two solutions of more general
hypergeometric equations maps the upper half-plane onto a curvilinear
triangle.

 Group-theoretic aspects of monodromy were studied by Jordan, and the
 geometrization of complex analysis around the concept of automorphic
 function was carried on by Klein, who once described his work as
 `blending Galois with Riemann'\footnote{quoted from [{\bf Gray86},
 p. 179].}.  The achievement fell to Poincar\'e. Starting from Fuchs'
 problem of finding all second-order differential equations for which a
 quotient of solutions $\tau$ admits a single-valued inverse, Poincar\'e
 founded the theory of `uniformizing differential equations' (in modern
 terminology), and eventually recognized that every Riemann surface of
 genus $>1$ should arise from the action of a discrete group of
 non-euclidean moves on the upper half-plane.

\bigskip
\bigskip

 What about non-archimedean analogues of this saga?

\bigskip

It is clear that the development of $p$-adic analysis and geometry took
a strikingly different route.


The very beginnings looked similar, indeed: J. Tate introduced rigid
geometry as a proper framework for the uniformization $\C_p^\times/q^\Z$
of elliptic curves with multiplicative reduction; B. Dwork developed
$p$-adic analysis starting from the hypergeometric differential equation
with parameters $(\frac{1}{2},\frac{1}{2};1)$, where he discovered the
essential notion of `Frobenius structure'. But it was immediately clear
that since solutions of $p$-adic differential equations do not converge
up to the next singularity, no faithful counterpart of complex monodromy
could take place.

From there on, different $p$-adic theories grew apart, with their own
languages, most of them relying on sophisticated parts of contemporary
algebraic geometry, and all claiming some analogy with ``the complex
case'':

\begin{itemize}
 \renewcommand{\labelitemi}{\normalfont\bfseries\textendash}
 \item the theory of differential equations matured slowly (with a strong
       orientation toward applications to exponential sums), struggling with
       the problems of singularities, without being able to tackle global
       problems,
 \item crystalline theory offered a global viewpoint on differential
       equations (oriented toward the cohomology of varieties in positive
       characteristic), but did not help to understand singularities,
 \item the theory of $p$-adic representations and $p$-adic Hodge theory
       developed independently of differential equations, as did the several
       avatars of rigid geometry, motivated by idiosyncratic problems.
\end{itemize}


 \bigskip

However, it is the author's opinion that the situation has somewhat
changed over the last years, that isolated branches are merging by fits
and starts\footnote{let us mention notably the incursion of rigid
geometry into the crystalline viewpoint and into the geometric theory of
finite coverings, the maturity of index theory and its applications to
algebraic geometry in positive characteristic, the new connections
between $p$-adic representations and differential equations on
annuli...}, and that a synthesis is gradually emerging. This encourages
to hope that, after many twists and turns, period mappings can indeed
become a unifying theme in the $p$-adic context. This book is intended
to be a contribution in this direction, by bringing closer the $p$-adic
theory to its complex precursor --- from periods and monodromy up to
triangle groups. Much remains to be done in order to achieve comparable
harmony and clarity.

\bigskip

\subsection*{Contents of this book.}

{\it Chapter I} is preparatory. It deals with the problems of analytic
continuation --- with emphasis on the case of solutions of differential
equations --- and periods of abelian integrals, in the $p$-adic context.

\medskip Multivalued complex-analytic functions can be handled in two
different, but essentially 
equivalent, ways:

$1)$ in a geometric way using Riemann surfaces, coverings and paths; 

$2)$ as limits of algebraic functions; this less orthodox way leads to a more
Galois-theoretic viewpoint on analytic continuation.

\medskip  Both ways are practicable in the $p$-adic context, but
eventually lead to completely different theories of analytic continuation. 

Following the first way demands to have at disposal $p$-adic spaces
which are locally arcwise connected.  Surprisingly enough in view of the
fractal nature of $p$-adic numbers, such a nice $p$-adic geometry does 
exist: it has been built by V. Berkovich (from the categorical
viewpoint, it is essentially equivalent to rigid geometry, in the sense
of Tate and Raynaud). In this framework, the monodromy of differential
equations can be analyzed in the usual way; but a new phenomenon occurs:
it is no longer true that there exists a basis of solutions around every
ordinary point. This approach is therefore limited to a rather special 
class of connections. We shall see in the sequel how the theory of
$p$-adic period mappings provides interesting examples in this class.
 
Following way $2)$, one encounters Dwork's notion of Frobenius
structure, which has often been considered as a plausible substitute for
monodromy in the $p$-adic context. We discuss so-called unit-root
crystals and overconvergence, and illustrate these notions by Dwork's
treatment of the $p$-adic Gamma function $\Gamma_p$, the Gross-Koblitz
formula, and by a detailed study of the $p$-adic hypergeometric function
$F(\frac{1}{2},\frac{1}{2},1; z)$.

\medskip

Abelian periods show themselves in two different ways:

$1)$ as integrals of algebraic differentials over loops, 

$2)$ when the integrand depends on a parameter $z$, as solutions of
certain linear differential equations in $z$ (Picard-Fuchs, or Gauss-Manin).

\medskip

Both ways are passable in the $p$-adic context, and again lead
to completely different theories. 

Following up the first way --- integrals being interpreted as Riemann sums
--- leads to P. Colmez' construction of abelian $p$-adic periods (in the
sense of Fontaine-Messing). These periods relate the Tate module 
($p$-primary torsion points) to the first De Rham cohomology module of
a given abelian variety over a $p$-adic field. They do not live in
$\C_p$, but in a certain $p$-adic ring $\matheur{B}^+_{\mathrm{dR}}$ (of which
$\C_p$ is a quotient), and it is not possible in general to express them 
function-theoretically, when the abelian variety moves in a
family. Nevertheless, we stress some particular cases where this can
actually be done, throwing a bridge between ways $1)$ and $2)$. 

The second way --- Dwork's viewpoint of periods as solutions of
Gauss-Manin connections --- will prevail at the modular level, when
$p$-divisible groups and filtered De Rham modules attached to families
of abelian varieties will be compared not directly, but via their moduli
spaces by the period mapping.

This way of looking at periods lacks arithmetic structure: namely, a
rational, or at least a $\ovl{\Q}$-structure, on the $\C_p$-space of
solutions, to convey some arithmetical meaning to the periods. We
discuss some cases where such a canonical structure exists (`$p$-adic
Betti lattices'), notably the case of abelian varieties with
supersingular reduction. We get here well-defined $p$-adic periods
(completely different from the Fontaine-Messing periods), and compute
them in terms of $\Gamma_p$-values in the case of elliptic curves with
complex multiplication ($p$-adic analogue of the Lerch-Chowla-Selberg
formula).


\bigskip

{\it Chapter II} is an introduction to the geometric theory of $p$-adic
period mappings, in the sense of Drinfeld-Rapoport-Zink.

We have tried to keep prerequisites at a minimum, and to emphasize as
much as possible the analogies between the complex and $p$-adic
contexts. Basic definitions about $p$-divisible groups and crystals are
recalled, and the proof of some basic results is sketched.

\medskip

The theory of period mappings attaches to a family of algebraic
varieties its periods, viewed abstractly as a moving point in a suitable
grassmannian. Due to constraints of Riemann-type, the period mapping
actually takes its values in an open subset of the grassmannian, the
period domain, which is a symmetric domain. It is multivalued, the
ambiguity being described by the projective monodromy of the Gauss-Manin
connection.

The younger $p$-adic theory is far less advanced: at present, there is a
theory of $p$-adic period mappings only for $p$-divisible groups or
closely related geometric objects. For want of wide range, it has
nevertheless gained richness and depth.

\medskip

In the presentation of Drinfeld-Rapoport-Zink, one starts by
constructing a moduli space $\MM$ for $p$-divisible groups which are
(quasi-)isogenous to a given one in characteristic $p$. The {\it
$p$-adic period mapping} $\PP$ then relates this moduli space to a
suitable grassmannian which parametrizes the Hodge filtration in the
Dieudonn\'e module. There are again constraints of Riemann-type, which
force the period mapping to take its values in an open subset of the
grassmannian, the period domain, which is a `symmetric domain'. The best
known example of such a period domain is the Drinfeld space
$\C_p\setminus \Q_p$, a $p$-adic analogue of the double-half-plane
$\C\setminus \R$.

In the situation where the $p$-divisible groups are algebraizable, \ie
come from $p$-primary torsion of abelian varieties parametri\-zed by a
certain Shimura variety $\Sh$, the moduli space $\MM$ provides a
uniformization of some tubular region in the Shimura variety $\Sh$. 

We show that, just as in the complex case, $\PP$ can be described in
terms of quotients of solutions of the associated Gauss-Manin connection
(this feature does not seem to appear in the literature, except in the
old special case of Dwork's period mapping for elliptic curves with 
ordinary reduction). This allows to give explicit formulas `\`a la
Dwork' for $\PP$ in many cases.

The most interesting cases, investigated by Drinfeld-Rapoport-Zink,
arise when $\Sh$ itself (more accurately: a whole connected component) is
uniformized by $\MM$: $\Sh$ is then a quotient of $\MM$ by an arithmetic
discrete group $\Gamma$.  We show that, up to replacing $\Sh$ by a finite
covering (to kill torsion in $\Gamma$), the solutions of the Gauss-Manin 
connection extend to global multivalued functions: in other words, the
Gauss-Manin connection has global monodromy in the 
sense of chapter I, and the projective monodromy group coincides with
$\Gamma$. 

We review the case of Shimura curves attached to quaternion algebras
over totally real fields (\v{C}erednik-Drinfeld-Boutot-Zink):
there is a global $p$-adic uniformization when $p$ divides the
discriminant of the quaternion algebra. We then construct, using the
$p$-adic Betti lattices of chapter I, a canonical $\ovl{\Q}$-space of
solutions of the Gauss-Manin connection which is stable under global
monodromy.



 \bigskip {\it Chapter III} explores the group-theoretic aspects of the
 theory of period mappings, in the $p$-adic context.

 The modern presentation of ramified coverings, uniformization, and of
 the Gauss-Riemann-Fuchs-Schwarz theory uses the notions of orbifold and
 uniformizing differential equation. We develop $p$-adic counterparts of
 these notions.

\medskip

The right notion of ramified covering to adopt is not obvious, because
\'etale coverings are only exceptionally topological coverings, in the
$p$-adic context. Ber\-ko\-vich's geometry provides a very convenient
framework for this kind of problems, since his spaces are locally arcwise
connected.  In the first section, we discuss the formalism of
fundamental groups attached to categories of (possibly infinite) \'etale
coverings satisfying simple axioms. Although the definition of these
topological groups is simple and natural, their topology itself may be
very complicated (not locally compact in general).

It turns out that there are too few topological coverings and too many
\'etale coverings, in general, to provide a workable theory of
monodromy. To remedy this, we introduce in the second section the notion
of {\it temperate \'etale covering}. Such coverings are essentially
built from (possibly infinite) topological coverings of finite \'etale
coverings. They are classified by temperate fundamental groups, which
seem to be the right $p$-adic equivalents, in many ways, of fundamental
groups of complex manifolds. These groups are not discrete in general,
but nevertheless often possess many infinite discrete quotients. We
present a large selection of examples.

\medskip Temperate \'etale coverings are well-suited to the definition
of $p$-adic orbifold charts and of the category of {\it $p$-adic
orbifolds} (they replace the unramified coverings, over $\C$). As in the
complex case, orbifolds and orbifold fundamental groups are the tools
for a theory of ramified coverings, developed in section 4.

Before turning to the $p$-adic analogue of the {\it uniformizing
differential equation} attached to an orbifold of dimension one (via
Schwarzian derivatives), we discuss local and global monodromy of
$p$-adic differential equations in section 3.  
 
We outline the Christol-Mebkhout theory of $p$-adic slopes of
differential equations over annuli, and the relation with Galois
representations of local fields of characteristic $p$.

We then define and study the {\it $p$-adic \'etale Riemann-Hilbert
functor}, which attaches a vector bundle with integrable connection to
any discrete representation of the \'etale fundamental group of a
$p$-adic manifold; connections in the image are characterized by the
fact that the \'etale sheaf of germs of solutions is locally
constant. This is a vast generalization of the phenomena of global
monodromy studied in Chapter I; for instance, the differential equation
$y'=y$ over the affine line belong to this class.

Uniformizing differential equations of orbifolds of dimension one also
belong to this class, and the representation actually factors through
the temperate fundamental group.  The case of a Shimura orbifold is of
special interest: the period mapping (complex or $p$-adic) is given by a
quotient of two solutions of a fuchsian differential equation defined
over a number field, which can be interpreted as uniformizing
differential equation either over $\C$ or over $\C_p$.

In Section 5, we examine the case of Schwarz orbifolds: the projective
line with $0,1,\infty$ as branched points (endowed with suitable
finite multiplicities). Over $\C$, the uniformizing differential
equations are of hypergeometric type, with projective monodromy group
identified with the orbifold fundamental group, namely with a
(cocompact) triangle group.

In the $p$-adic case, we define {\it $p$-adic triangle groups} to be
projective monodromy groups of those hypergeometric differential
equations which are in the image of the $p$-adic \'etale Riemann-Hilbert
functor: thus by definition, there exists a finite \'etale covering of
$\C_p\setminus \{0,1\}$ over which the hypergeometric function extends
to a global multivalued analytic function.

We give a purely geometric description of these discrete subgroups of
$PSL_2(\C_p)$ (without reference to differential equations; it is
perhaps here that we are closest to Fuchs and Schwarz).  From this, and
recent combinatorial work by F. Kato, it follows that infinite $p$-adic
triangle groups exist only for $p\leq 5$.

We then construct the $p$-adic analogues of Takeuchi's list of
arithmetic triangle groups, \ie the list of all ``arithmetic'' $p$-adic
triangle groups for every $p$, using the 
\v{C}erednik-Drinfeld-Boutot-Zink uniformization of Shimura curves and
$p$-adic period mapping. Special values of the corresponding
hypergeometric functions at CM points are expressed in terms of
$\Gamma_p$-values.

\medskip

\subsection*{On the style.} 

\medskip

The first chapter is rather down-to-earth. It has kept
something of the informal style of lecture notes, and this also holds to
a lesser extent for the second chapter; the level is inhomogeneous, the
pace may sometimes be brisk, proofs are often omitted and replaced by
references.  In contrast, the last and longest chapter is devoted to a
systematic exposition of new material.

We hope that our constant function-theoretic viewpoint brings some unity
to the exposition, throughout the chapters.  We have tried to keep them
(and even the sections) as logically independent as possible. Thus, for
example:

only subsections 1.5 and 5.3 of chapter I are needed in the sequel;

until section 7, chapter II is almost self-contained;

until subsection 4.7, chapter III is almost self-contained;

reading fragments of III should be enough to grasp the ins and outs
of $p$-adic triangle groups.
 
\medskip

Sections I.3 and III.6 are small pieces of `computational mathematics'
(without computer) intended as testing ground for notions developed in
the preceding sections.

\bigskip

{\bf Acknowledgments.}
It is a pleasure to acknowledge here our main sources of inspiration:


first of all, V. Berkovich's views on $p$-adic geometry: without this
extremely convenient way of seeing $p$-adic spaces, our systematic
to-ings and fro-ings between complex and $p$-adic worlds would have 
been impossible;

Rapoport-Zink's book on period spaces for $p$-divisible groups has been
a constant reference for Chapter II;

in Chapter III, our treatment has been much influenced by the work of
M. van der Put on $p$-adic discontinuous groups, by J. De Jong's study
of \'etale coverings, and no less by M. Yoshida's orbifold viewpoint on
the (complex) Gauss-Schwarz theory.

\medskip

Concerning this book, I am much indebted to F. Kato in many ways: not
only did he organize the T\={o}hoku and the Hokkaid\={o} lectures
series, asking me to speak on these matters, but he also suggested the
very project of writing these notes, and influenced a lot both the
content and the style by his numerous questions and comments; I have
been much encouraged in writing them by his constant interest and by his
own recent work on $p$-adic triangle groups.

\medskip

The T\={o}hoku lectures, as well as the birth of some ideas presented in
III (especially $p$-adic triangle groups), took place during my stay in
Japan thanks to a fellowship from the Japan Society for the promotion of
Science. I am very grateful to this institution, and to H. Shiga who
arranged the stay.

\medskip

I thank M. Matignon, F. Baldassarri and S. Bosch for invitations to
lecture on these topics at the Universities of Bordeaux, Padova, and
M\"unster respectively.

 \medskip

I am grateful to F. Kato and N. Tsuzuki for writing the appendices of
this book. I am also very grateful to S. Matsuda who kindly took care,
with great patience, of the nightmarish electronic harmonization of the
main text and the appendices, and drawing the figures.

 \medskip

Thanks are also due to P. Colmez and O. Gabber for several useful
comments, and to J.F. Boutot for sending me overseas his joint preprint
with T. Zink about uniformization of Shimura curves.

I also thank very much P. Bradley for his careful reading of a version
of chapter III and for kindly providing a long list of errata (of
course, I am responsible for any remaining error).
